\vspace{-3ex}
\section{Introduction}
Vision-language pre-training has recently received tremendous success on various multimodal downstream tasks.
However, existing methods have two major limitations:

(1) Model perspective: most methods either adopt an encoder-based model~\cite{clip,ALBEF}, or an encoder-decoder~\cite{VL_T5,simvlm} model.
However, encoder-based models are less straightforward to directly transfer to text generation tasks (\eg~image captioning), whereas encoder-decoder models have not been successfully adopted for image-text retrieval tasks.

(2) Data perspective: most state-of-the-art methods (\eg, CLIP~\cite{clip}, ALBEF~\cite{ALBEF}, SimVLM~\cite{simvlm}) pre-train on image-text pairs collected from the web. 
Despite the performance gain obtained by scaling up the dataset,
our paper shows that the noisy web text is suboptimal for vision-language learning.


To this end, we propose \name: Bootstrapping Language-Image Pre-training for unified vision-language understanding and generation.
\name~is a new VLP framework which enables a wider range of downstream tasks than existing methods.
It introduces two contributions from the model and data perspective, respectively:

(a) Multimodal mixture of Encoder-Decoder (MED):
a new model architecture for effective multi-task pre-training and flexible transfer learning.
An MED can operate either as a unimodal encoder, or an image-grounded text encoder, or an image-grounded text decoder.
% The image features are passed in through cross-attention layers.
% To le multi-task pre-training,
% the MED model reuses its parameters across the encoder-mode and decoder-mode, except for the self-attention layers.
% The text encoder and text decoder share the same transformer parameters except for the self-attention layers.
The model is jointly pre-trained with three vision-language objectives: image-text contrastive learning, image-text matching, and image-conditioned language modeling.
% We propose a parameter-sharing strategy to leverage multi-task learning with efficiency.


% For the purpose of BootCap,
% we finetune the image-grounded text decoder into the \textit{generator},
% and the image-grounded text encoder into the \textit{denoiser}.

(b) Captioning and Filtering (CapFilt):
a new dataset boostrapping method for learning from noisy image-text pairs.
We finetune a pre-trained MED into two modules: a \textit{captioner} to produce synthetic captions given web images,
and a \textit{filter} to remove noisy captions from both the original web texts and the synthetic texts.
% We observe substantial performance gains by injecting diverse synthetic captions into the pre-training corpus. 

\begin{figure}
\vspace{-1ex}
    \centering 
\includegraphics[width=1\columnwidth]{imgs/teaser.pdf}
    \vspace{-5.5ex}
    \caption{We use a Captioner (Cap) to generate synthetic captions for web images, and a Filter (Filt) to remove noisy captions.}
    \label{fig:teaser}
    \vspace{-2.5ex}
\end{figure}

We perform extensive experiments and analysis, and make the following key observations. 
\vspace{-\topsep}
\vspace{-1ex}
\begin{itemize}[leftmargin=*]
	\setlength\itemsep{0pt}
% 	\item Our pre-trained model transfers flexibly to both understanding-based and generation-based tasks, in the meantime being efficient during pre-training.	
	\item We show that the captioner and the filter work together to achieve substantial performance improvement on various downstream tasks by bootstrapping the captions.  
	We also find that more diverse captions yield larger gains.
	\item \name~achieves state-of-the-art performance on a wide range of vision-language tasks,
	including image-text retrieval, image captioning, visual question answering, visual reasoning, and visual dialog. We also achieve state-of-the-art zero-shot performance when directly transferring our models to two video-language tasks: text-to-video retrieval and videoQA.
\end{itemize}